\documentclass[11pt]{article}
\usepackage[letterpaper,margin=0.8in]{geometry}
\usepackage[utf8]{inputenc}
\usepackage[T1]{fontenc}
\usepackage{amsmath,amssymb,mathtools,amsthm}
\usepackage{newpxtext,newpxmath}
\usepackage{booktabs,microtype}
\usepackage[hidelinks]{hyperref}
\hypersetup{
  pdftitle={The Gauge Algebra of Chirality},
  pdfauthor={Emad Mostaque},
  pdfsubject={An ordered two-simple-factor adjoint-complement classification}
}
\newtheorem{theorem}{Theorem}
\newtheorem{corollary}[theorem]{Corollary}
\newcommand{\Tr}{\operatorname{Tr}}
\newcommand{\ad}{\operatorname{ad}}
\newcommand{\su}{\mathfrak{su}}
\newcommand{\so}{\mathfrak{so}}
\newcommand{\spalg}{\mathfrak{sp}}
\newcommand{\sectionhead}[1]{%
  \par\addvspace{7pt}{\large\bfseries #1}\par\nobreak\vspace{2pt}}
\setlength{\parindent}{0pt}
\setlength{\parskip}{4pt plus .5pt minus .5pt}
\setlength{\abovedisplayskip}{6pt}
\setlength{\belowdisplayskip}{6pt}
\setlength{\abovedisplayshortskip}{3pt}
\setlength{\belowdisplayshortskip}{3pt}
\emergencystretch=2em
\pagestyle{plain}

\begin{document}
\begin{center}
{\LARGE\bfseries The Gauge Algebra of Chirality\par}
\vspace{5pt}
Emad Mostaque\par
April 2026
\end{center}
\vspace{1pt}
\hrule
\vspace{6pt}

We classify adjoint complements of ordered maximal two-simple-factor
subalgebras of compact simple Lie algebras by complex type and cubic trace.

\sectionhead{The domain and the two tests}
Let \(\mathfrak g\) be compact and simple, and let
\(\mathfrak h=\mathfrak k\oplus\mathfrak f\subsetneq\mathfrak g\) have
\emph{exactly two nonabelian simple ideals}. Require \(\mathfrak h\) to be
maximal among all proper Lie subalgebras, not merely such products.
Both orderings are tested, with \(\mathfrak k\) the tested factor.
The invariant orthogonal complement
\(\mathfrak q=\mathfrak h^\perp\) is an \(\mathfrak h\)-module.

Decompose \(\mathfrak q_{\mathbb C}\) into \emph{joint}
\(\mathfrak h\)-irreducibles \(R\boxtimes S\).
Put all summands with \(R\simeq R^*\) in \(q_0\).
The rest occur in dual pairs with equal multiplicities:
\[
 \mathfrak q_{\mathbb C}
 =q_0\oplus\bigoplus_\alpha m_\alpha(U_\alpha\oplus U_\alpha^*).
\]
Choose one member of each such pair, retaining its \emph{full multiplicity},
and call their sum \(W\). Every choice is checked. Thus \(W\) is
\(\mathfrak h\)-invariant: one may not independently reorient the
\(\mathfrak k\)-copies inside an \(\mathfrak f\)-multiplet.
The tests are
\[
 \text{(A)}\quad W\ne0,
 \qquad
 \text{(B)}\quad
 d_{\mathfrak k}(W)(X,Y,Z)
 =\tfrac12\Tr_W\!\left(T_X\{T_Y,T_Z\}\right)=0
 \quad(X,Y,Z\in\mathfrak k),
\]
where \(T_X=-i\rho_W(X)\) is Hermitian. Test (B) is the vanishing of the
local cubic anomaly tensor on \(\mathfrak k\).

\sectionhead{Exclusion at every classical rank}
The classification [1,2] gives parameterized families of unbounded rank.
Unitary block stabilizers have a central \(\mathfrak u(1)\), outside this
domain. Orthogonal and symplectic blocks have tensor-product defining
complements with self-dual constituents. The tensor-embedding families are
\[
\begin{array}{ll}
 \su(p)\oplus\su(q)\subset\su(pq),&
 \so(p)\oplus\so(q)\subset\so(pq),\\
 \spalg(p)\oplus\spalg(q)\subset\so(4pq),&
 \so(p)\oplus\spalg(q)\subset\spalg(pq).
\end{array}
\]
Here \(\spalg(p)\) has defining dimension \(2p\). Retain only maximal
inclusions satisfying the simplicity requirements, including low-rank
isomorphisms. The unitary tensor complement is
\(\ad\su(p)\boxtimes\ad\su(q)\). In the other cases use
\[
 \Lambda^2(U\otimes Z)=
 (\Lambda^2U\otimes\mathrm{Sym}^2Z)
 \oplus(\mathrm{Sym}^2U\otimes\Lambda^2Z)
\]
and the analogous formula for \(\mathrm{Sym}^2(U\otimes Z)\), with both
symmetric squares together and both exterior squares together.
After subtracting the two adjoints, the factorwise constituents are
self-dual: adjoints and traceless symmetric or exterior squares of defining
modules. Hence every classical row has \(W=0\), at arbitrary rank.

\sectionhead{Exceptional elimination}
The exceptional branchings [3] leave the five unordered embeddings on
page~2. In particular, for \(E_6\supset A_2+G_2\), with factors in that order,
\[
 \mathbf{78}=(\mathbf8,\mathbf1)\oplus(\mathbf1,\mathbf{14})
                    \oplus(\mathbf8,\mathbf7).
\]
Both factors see only self-dual constituents. The complex \((\mathbf3,\mathbf7)\)
belongs instead to
\(\mathbf{27}=(\overline{\mathbf6},\mathbf1)\oplus(\mathbf3,\mathbf7)\).
The \(A_5+A_1\) complement \((\mathbf{20},\mathbf2)\) is also factorwise
self-dual. Neither embedding passes (A).

\newpage
\sectionhead{The five complex rows and the survivor}
Dynkin types label the factors below. Entries are cubic coefficients in
the convention \(A(\mathbf n)=1\) for \(\su(n)\); dualizing a whole half
reverses their signs. The reverse column also reapplies test (A).
\begin{center}
\renewcommand{\arraystretch}{1.16}
\setlength{\tabcolsep}{7pt}
\begin{tabular}{@{}lllll@{}}
\toprule
\(\mathfrak g\)&\((\mathfrak k,\mathfrak f)\)&Half \(W\)&
\(\mathfrak k\) test&Reverse test\\
\midrule
\(F_4\)&\((A_2,A_2)\)&\((\mathbf6,\overline{\mathbf3})\)&\(21\)&\(-6\)\\
\(E_7\)&\((A_5,A_2)\)&\((\mathbf{15},\overline{\mathbf3})\)&\(6\)&\(-15\)\\
\(E_8\)&\((A_4,A_4)\)&\(W_{a,b}\)&\(5a+10b\)&\(-10a+5b\)\\
\(E_8\)&\((A_2,A_1)\)&\((\mathbf{10},\mathbf5)\)&\(135\)&(A) fails\\
\(E_8\)&\((E_6,A_2)\)&\((\mathbf{27},\mathbf3)\)&\(0\)&\(27\)\\
\bottomrule
\end{tabular}
\end{center}
For the \(A_4+A_4\) row, put
\(U=(\mathbf{10},\overline{\mathbf5})\),
\(Z=(\mathbf5,\mathbf{10})\), and
\(W_{a,b}=U^{(a)}\oplus Z^{(b)}\), where \(a,b=\pm1\) and
\(R^{(+1)}=R,\ R^{(-1)}=R^*\). Both coefficient sets are
\(\{\pm5,\pm15\}\). For the special \(A_2+A_1\subset E_8\) embedding,
the remaining complement is
\(q_0=(\mathbf8,\mathbf7)\oplus(\mathbf{27},\mathbf3)\), with self-dual
\(A_2\)-constituents; in reverse order every \(A_1\)-representation is self-dual.

The calculations use
\(A(R\boxtimes S)=\dim(S)A(R)\), additivity and duality, with
\(A_{A_2}(\mathbf6)=7\),
\(A_{A_5}(\mathbf{15})=2\), and
\(A_{A_2}(\mathbf{10})=27\).
The \(E_6\) coefficient vanishes because its Lie algebra has no invariant
symmetric cubic. This is distinct from the cubic on
\(\mathbf{27}^{\otimes3}\), which has representation rather than
Lie-algebra indices.

\begin{theorem}[Ordered maximal-product classification]
In the stated domain and with the stated half-selection rule,
(A) and (B) select
\(\mathfrak e_8\supset\mathfrak e_6\oplus\su(3)\),
with tested factor \(\mathfrak k=\mathfrak e_6\), uniquely up to
isomorphism of ordered Lie-algebra pairs. The half is
\(W=\mathbf{27}\boxtimes\mathbf3\), up to duality.
\end{theorem}

\sectionhead{Restriction and net chiral content}
The adjoint branching is
\[
 \mathbf{248}=(\mathbf{78},\mathbf1)\oplus(\mathbf1,\mathbf8)
 \oplus(\mathbf{27},\mathbf3)
 \oplus(\overline{\mathbf{27}},\overline{\mathbf3}).
\]
Consequently \(W|_{E_6}=3\,\mathbf{27}\). At group level use
\[
 G_{\mathrm{SM}}=
 [SU(3)\times SU(2)\times U(1)]/\mathbb Z_6
 \subset SU(5)\subset\operatorname{Spin}(10)\subset E_6^{\mathrm{sc}}.
\]
Restriction gives
\[
 \mathbf{27}|_{\operatorname{Spin}(10)}
 =\mathbf{16}\oplus\mathbf{10}\oplus\mathbf1,\qquad
 \mathbf{27}|_{SU(5)}
 =\underbrace{\mathbf{10}\oplus\overline{\mathbf5}}_{\mathcal F}
   \oplus(\mathbf5\oplus\overline{\mathbf5})\oplus2\,\mathbf1.
\]
Here \(\mathcal F|_{G_{\mathrm{SM}}}\) is one minimal family:
\[
 (\mathbf3,\mathbf2)_{1/6}\oplus
 (\overline{\mathbf3},\mathbf1)_{-2/3}\oplus
 (\overline{\mathbf3},\mathbf1)_{1/3}\oplus
 (\mathbf1,\mathbf2)_{-1/2}\oplus(\mathbf1,\mathbf1)_1.
\]
The labels use the all-left-handed convention; the singlet in the
\(\mathbf{16}\) is the additional \(\nu^c\), not part of this minimal family.

\begin{corollary}[Containment]
The selected half restricts to three minimal Standard Model families,
three vectorlike \(\mathbf5\oplus\overline{\mathbf5}\) pairs, and six
singlets. Its net chiral class is
\(3\chi(\mathcal F|_{G_{\mathrm{SM}}})\); the dual half
has the opposite class.
\end{corollary}
Here \(\chi\) records representation-minus-dual multiplicities and
annihilates self-dual summands.

\par\addvspace{5pt}
{\small
\textbf{References.}
[1] E.\,B.~Dynkin, \emph{Semisimple subalgebras of semisimple Lie algebras},
Mat.\ Sb.\ \textbf{30} (1952), 349--462;
\emph{Maximal subgroups of the classical groups},
Tr.\ Mosk.\ Mat.\ Obs.\ \textbf{1} (1952), 39--166.
[2] F.~Antoneli, M.~Forger and P.~Gaviria,
\emph{Maximal Subgroups of Compact Lie Groups},
\href{https://arxiv.org/abs/math/0605784}{arXiv:math/0605784}, Tables 2--4.
[3] R.~Slansky, \emph{Group theory for unified model building},
Phys.\ Rep.\ \textbf{79} (1981), 1--128,
\href{https://doi.org/10.1016/0370-1573(81)90092-2}{doi:10.1016/0370-1573(81)90092-2},
Tables 15 and 58.
}
\end{document}
